\documentclass[10pt]{article}
\usepackage[T1]{fontenc}
\usepackage[utf8]{inputenc}
\usepackage[french]{babel}
\usepackage{amsmath,amssymb,mathrsfs,mathpazo}
\usepackage{geometry}
\geometry{a4paper,margin=2cm}
\usepackage{graphicx,array,tabularx,tabularray,colortbl,multirow,multicol,booktabs,enumitem,pifont,ulem,fancybox,diagbox,cancel}
\usepackage{pstricks,pstricks-add,pst-plot,pst-tree,pst-node,pst-text,pst-func,pst-math,pst-bspline,pst-3dplot,pst-eucl,pst-solides3d}
\let\pscalculate\undefined
\usepackage{pst-bezier}
\usepackage{tikz,pgfplots,tkz-tab}
\usetikzlibrary{arrows,arrows.meta,patterns,calc,positioning,trees,decorations.markings,shapes.geometric,shapes.misc,shapes.arrows,intersections,decorations.pathmorphing,decorations.pathreplacing,backgrounds,shadows,babel}
\pgfplotsset{compat=1.18}
\usepgfplotslibrary{fillbetween}
\usepackage[np]{numprint}
\usepackage{listings,makecell,pas-tableur,twoopt}
\let\barre\undefined
\usepackage[pstricks]{tablvar}
\definecolor{darkgreen}{RGB}{0,100,0}
\providecommand{\red}{\color{red}}
\providecommand{\blue}{\color{blue}}
\providecommand{\green}{\color{green}}
\providecommand{\black}{\color{black}}
\providecommand{\white}{\color{white}}
\providecommand{\gray}{\color{gray}}
\providecommand{\cyan}{\color{cyan}}
\providecommand{\magenta}{\color{magenta}}
\providecommand{\yellow}{\color{yellow}}
\providecommand{\violet}{\color{violet}}
\providecommand{\orange}{\color{orange}}
\providecommand{\euro}{\texteuro}
\providecommand{\EUR}{\texteuro}
\providecommand{\ieme}{\textsuperscript{e}}
\providecommand{\degres}{\ensuremath{^\circ}}
\providecommand{\degree}{\ensuremath{^\circ}}
\providecommand{\vv}[1]{\overrightarrow{#1}}
\providecommand{\Coord}[1]{#1}
\providecommand{\decosix}{}
\providecommand{\starredbullet}{\ensuremath{\bullet}}
\pagestyle{empty}
\setlength{\parindent}{0pt}
\usepackage{adjustbox}\DeclareUnicodeCharacter{2236}{:}\DeclareUnicodeCharacter{2190}{\ensuremath{\leftarrow}}
\providecommand{\R}{\mathbb{R}}
\providecommand{\N}{\mathbb{N}}
\providecommand{\Z}{\mathbb{Z}}
\providecommand{\Q}{\mathbb{Q}}
\providecommand{\D}{\mathbb{D}}
\providecommand{\C}{\mathbb{C}}
\providecommand{\e}{{\mathrm{e}}}
\providecommand{\vect}[1]{\overrightarrow{#1}}
\providecommand{\vectt}[1]{\overrightarrow{\mathrm{#1}}}
\providecommand{\vvt}[1]{\overrightarrow{\mathrm{#1}}}
\providecommand{\np}[2][]{{#2~#1}}
\providecommand{\vv}[1]{{\overrightarrow{#1}}}
\renewcommand{\d}{{\mathrm{d}}}
\providecommand{\degres}{{^\circ}}
\providecommand{\degre}{{^\circ}}
\providecommand{\blue}{\color{blue}}
\providecommand{\red}{\color{red}}
\providecommand{\enskip}{\quad}
\providecommand{\llbracket}{\lbrack\!\lbrack}
\providecommand{\rrbracket}{\rbrack\!\rbrack}
\providecommand{\euro}{{\text{€}}}
\providecommand{\no}{{\text{n°}}}
\providecommand{\interval}[2]{{[#1;#2]}}
\providecommand{\intervFO}[2]{{[#1;#2[}}
\providecommand{\dotfill}{\ldots}

\providecommand{\ds}{\displaystyle}
\providecommand{\textding}[1]{#1}
\providecommand{\Oij}{\(\mathrm{O};\vec{\imath},\vec{\jmath}\)}
\providecommand{\Oijk}{\(\mathrm{O};\vec{\imath},\vec{\jmath},\vec{k}\)}
\providecommand{\Ouv}{\(\mathrm{O};\vec{u},\vec{v}\)}
\providecommand{\intervFF}[2]{\ensuremath{[#1;#2]}}
\renewcommand{\euro}{\ifmmode\text{\texteuro}\else\texteuro\fi}
\DeclareUnicodeCharacter{00B0}{\ifmmode^{\circ}\else\textdegree\fi}
\DeclareUnicodeCharacter{20AC}{\ifmmode\text{\texteuro}\else\texteuro\fi}
% Énoncés : Éducation nationale. Transcription et corrigés : APMEP, Denis Vergès et contributeurs. Découpage et classement : CoopMaths, Sébastien Lozano, Éric Elter et contributeurs. Adaptation de la présentation : Maths Explorer.
% CC BY-SA : https://coopmaths.fr/www/about/
% https://coopmaths.fr/alea/static/bac/2023/tex/bac_2023_03_sujet1_polynesie_2_cor.tex
\begin{document}
\noindent\textit{APMEP / CoopMaths. CC BY-SA. Présentation adaptée par Maths Explorer.}\par\medskip
\medskip

\begin{enumerate}
\item 
	\begin{enumerate}
		\item L'équation paramétrique de $d_2$ montre qu'elle contient le point de coordonnées $(-3~;~0~;~5)$ et de vecteur directeur dont les composantes sont les coefficients de $k$ donc $\vect{v}\begin{pmatrix}2\\1\\0\end{pmatrix}$.
		\item Les vecteurs directeurs $\vect{u}$ et $\vect{v}$ ne sont manifestement pas colinéaires, donc les droites $d_1$ et $d_2$ ne sont pas parallèles.
		\item Une représentation paramétrique de la droite $d_1$ est $\left\{\begin{array}{l c l}
x&=2 + t\\
y&=&3 - t\\
z&=&t
\end{array}\right. \quad , t \in \R$.

S'il existe un point commun aux deux droites, il doit donc exister des réels $k$ et $t$ tels que :

$\left\{\begin{array}{l c l c l}
x &=& 2 + t&=&2k - 3\\
y&=&3 - t&=&k\\
z&=&t&=&5
\end{array}\right.$

La dernière équation donne $t = 5$, puis la deuxième $k = 3 - 5 = - 2$ et en remplaçant dans la première $2 + 5 = - 4 - 3$ : cette égalité est fausse, donc $d_1$ et $d_2$  ne sont pas sécantes.
		\item Les droites n'étant ni sécantes ni parallèles elles ne sont pas coplanaires.
	\end{enumerate}
\item 
	\begin{enumerate}
		\item On a $\vect{w}\ \cdot \vect{u} = - 1 - 2 + 3 = 0$ ;
		
		de même $\vect{w}\ \cdot \vect{v} = - 2 + 2 + 0 = 0$, donc le vecteur $\vect{w}$ est orthogonal aux vecteurs $\vect{u}$ et $\vect{v}$.
		\item Un point commun au plan $P$ et à la droite $d_2$ a ses coordonnées qui vérifient le système :
		
$\left\{\begin{array}{l c l}
x &=&2k - 3\\
y &=&k\\
z &=&5\\
5x + 4y - z - 22 &=& 0
\end{array}\right. \quad , t \in \R$, d'où en remplaçant $x,\:y,\:z$ dans la dernière équation :

$5(2k - 3) + 4(k) - 5 - 22 = 0\iff 10k - 15 + 4k - 27 = 0\iff 14k = 42 \iff 7k = 21 \iff k = 3$.

Le point commun au plan $P$ et à la droite $d_2$ a pour coordonnées (6 - 3~;~3~;~5) soit M(3~;~3~;~5).
	\end{enumerate}
\item 	
	\begin{enumerate}
		\item $\Delta$ et $d_1$ ont leurs vecteurs directeurs $\vect{w}$ et $\vect{u}$  orthogonaux donc ces droites sont orthogonales.
		
Ces deux droites sont sécantes s'il existent des réels $t$ et $r$ tels que 

$\left\{\begin{array}{l c l c l}
x&=&2 + t&=&- r + 3\\
y&=&3 - t&=&2r + 3\\
z&=&t   &=&3r + 5
\end{array}\right. \quad , t, r  \in \R$.

La dernière équation $t = 3r + 5$ donne en remplaçant dans la deuxième : 

$3 - 3r - 5 = 2r + 3 \iff - 5 = 5r \iff r = - 1$, d'où $t = 5 - 3 = 2$ et en remplaçant dans la première équation on obtient $2 + 2 = 1 + 3$, égalité vraie.

$\Delta$ et $d_1$ sont sécantes au point L(4~;~1~;~2)
		\item  Conclusion : on a trouvé une droite $\Delta$ perpendiculaire commune aux deux droites $d_1$ et $d_2$.
	\end{enumerate}
\end{enumerate}

\bigskip

\end{document}
